\documentclass[12pt]{article}
%^ Start of document
\usepackage{times}  %Makes text TimesNewRoman
\usepackage{setspace} %Allows you to set single or double space 
\usepackage{biblatex} %Imports biblatex package
\usepackage{graphicx} % Required for inserting images
\usepackage{authblk}
\usepackage[colorlinks=true, urlcolor=blue, linkcolor=blue]{hyperref}
\usepackage{subfig}
\usepackage{float}
\usepackage{tabularx}
\usepackage{multirow}
\usepackage[paperheight=11in,
   paperwidth=8.5in,
   top=20mm,
   bottom=20mm,
   left=40mm,
   right=40mm]{geometry}
\usepackage{moresize}
\usepackage{tikz}
\usepackage{xcolor}
\usepackage{physics}
\usepackage{amssymb}
\usepackage{mathrsfs}
\usepackage{amsmath}
\usepackage{ragged2e}
\usepackage{comment}
\usepackage{xcolor}
\addbibresource{mathmeth.bib}
\begin{document}
\singlespacing  %Sets single spacing for whole document
\title{Math Methods HW 3}

\author[1]{Zachary Tullier}
\affil[1]{Student\\Physics Department\\ University of Houston - Clear Lake \\Houston, Texas\\U.S}
\date{}
\maketitle

\section*{Textbook Question 7.4.2}
    Show that Legendre's equation has regular singularities at $x=-1$, 1, and $\infty$

    \subsection*{Solution:}
        The standard form of Legendre’s equation is:
        \[
        (1 - x^2) \frac{d^2y}{dx^2} - 2x \frac{dy}{dx} + n(n+1)y = 0,
        \]
        where \(n\) is a constant.
        A point \(x = x_0\) is a regular singular point of a second-order differential equation if, after rewriting the equation in standard form:
        \[
        \frac{d^2y}{dx^2} + P(x)\frac{dy}{dx} + Q(x)y = 0,
        \]
        the functions \((x - x_0)P(x)\) and \((x - x_0)^2Q(x)\) are analytic (finite and well-defined) at \(x = x_0\). Points where the equation fails to be analytic are singular points, and if these conditions hold, they are regular singular points.
        The equation can be rewritten as:
        \[
        \frac{d^2y}{dx^2} - \frac{2x}{1 - x^2} \frac{dy}{dx} + \frac{n(n+1)}{1 - x^2} y = 0.
        \]
        Here:
        \[
        P(x) = -\frac{2x}{1 - x^2}, \quad Q(x) = \frac{n(n+1)}{1 - x^2}.
        \]
        The denominators \(1 - x^2 = 0\) give the singular points of the equation, which are at \(x = \pm 1\). Additionally, \(x = \infty\) needs to be examined separately. At \(x = 1\):
        \begin{enumerate}
            \item \((x - 1)P(x)\)
            \[
            P(x) = -\frac{2x}{1 - x^2} = -\frac{2x}{(1 - x)(1 + x)}.
            \]
            Near \(x = 1\), \(P(x) \approx -\frac{2}{1 + x} \cdot \frac{1}{x - 1}\), so:
            \[
            (x - 1)P(x) = -\frac{2}{1 + x},
            \]
            which is analytic at \(x = 1\).
            \item \((x - 1)^2Q(x)\)
            \[
            Q(x) = \frac{n(n+1)}{1 - x^2} = \frac{n(n+1)}{(1 - x)(1 + x)}.
            \]
            Near \(x = 1\), \(Q(x) \approx \frac{n(n+1)}{(x - 1)(1 + x)}\), so:
            \[
            (x - 1)^2Q(x) = \frac{n(n+1)(x - 1)}{1 + x}.
            \]
            This is analytic at \(x = 1\).
        \end{enumerate}
        Thus, \(x = 1\) is a regular singular point. At \(x = -1\) the process is similar to \(x = 1\). Rewrite \(P(x)\) and \(Q(x)\) and check the terms:
        \[
        (x + 1)P(x) \quad \text{and} \quad (x + 1)^2Q(x)
        \]
        are both analytic at \(x = -1\).
        Thus, \(x = -1\) is a regular singular point.
        To analyze \(x = \infty\), use the substitution \(z = \frac{1}{x}\), so that \(x \to \infty\) corresponds to \(z \to 0\). Rewrite Legendre’s equation in terms of \(z\):
        \[
        (1 - z^2) \frac{d^2y}{dz^2} + 2z \frac{dy}{dz} + n(n+1)y = 0.
        \]
        The singularity at \(z = 0\) corresponds to \(x = \infty\). Following a similar analysis as above, the terms \(zP(z)\) and \(z^2Q(z)\) are analytic at \(z = 0\).
        Thus, \(x = \infty\) is a regular singular point.

\section*{Textbook Question 7.5.5}
    Obtain a series solution of the hypergeometric equation
    \[
        x(x-1)y'' + \left((1 + a + b)x - c \right) y' + aby = 0
    \]
   
   \subsection*{Solution:}
        Assume the solution \(y(x)\) is of the form:
        \[
        y(x) = \sum_{n=0}^\infty a_n x^{n+s},
        \]
        where \(s\) is a parameter to be determined, and \(a_n\) are the coefficients.
        The derivatives of \(y(x)\) are:
        \[
        y'(x) = \sum_{n=0}^\infty a_n (n+s)x^{n+s-1},
        \]
        \[
        y''(x) = \sum_{n=0}^\infty a_n (n+s)(n+s-1)x^{n+s-2}.
        \]
        Substitute \(y(x)\), \(y'(x)\), and \(y''(x)\) into the hypergeometric equation:
        \[
        x(x-1)y'' + [(1+a+b)x - c]y' + aby = 0.
        \]
        Each term becomes:
        \begin{enumerate}
            \item For \(x(x-1)y''\):
            \[
            x(x-1)y'' = x(x-1)\sum_{n=0}^\infty a_n (n+s)(n+s-1)x^{n+s-2}.
            \]
            Simplify:
            \[
            x(x-1)y'' = \sum_{n=0}^\infty a_n (n+s)(n+s-1)x^{n+s} - \sum_{n=0}^\infty a_n (n+s)(n+s-1)x^{n+s+1}.
            \]
        
            \item For \([(1+a+b)x - c]y'\):
            \[
            [(1+a+b)x - c]y' = \sum_{n=0}^\infty a_n (n+s)(1+a+b)x^{n+s} - \sum_{n=0}^\infty a_n (n+s)c x^{n+s-1}.
            \]
            \item For \(aby\):
            \[
            aby = \sum_{n=0}^\infty ab a_n x^{n+s}.
            \]
        \end{enumerate}
        Combine all terms and equate coefficients of \(x^{n+s}\) to zero.
        The indicial equation arises from the lowest power of \(x\), which occurs when \(n = 0\). Collecting terms proportional to \(x^s\) gives:
        \[
        s(s-1) + s(1+a+b-c) + ab = 0.
        \]
        Simplify:
        \[
        s^2 + (a+b-c)s + ab = 0.
        \]
        This is a quadratic equation for \(s\), whose roots are:
        \[
        s = \frac{c - (a+b)}{2} \quad \text{and} \quad s = \frac{c - (a+b) + 2}{2}.
        \]
        For \(n \geq 1\), collect terms proportional to \(x^{n+s}\):
        \[
        a_n \left[(n+s)(n+s-1) + (n+s)(1+a+b-c) + ab\right] + a_{n-1}(n+s-1)(1+a+b) - a_{n-1}c = 0.
        \]
        This simplifies to the recurrence relation:
        \[
        a_n = \frac{(n+s-1)(c - (1+a+b))a_{n-1}}{(n+s)(n+s-1) + (n+s)(1+a+b-c) + ab}.
        \]
        The general solution is then written as:
        \[
        y(x) = \sum_{n=0}^\infty a_n x^{n+s},
        \]
        where the coefficients \(a_n\) are determined recursively using the relation above, with the initial condition \(a_0 \neq 0\) (arbitrary constant).
        To test for convergence, we apply the ratio test:
        \[
        \lim_{n \to \infty} \left| \frac{a_{n+1}x^{n+1+s}}{a_nx^{n+s}} \right| = \lim_{n \to \infty} \left| \frac{a_{n+1}}{a_n} \cdot x \right|.
        \]
        From the recurrence relation, for large \(n\), the coefficients \(a_n\) behave as:
        \[
        a_n \sim \frac{1}{n!}.
        \]
        Thus, the series converges for \(|x| < 1\), which is the region of convergence for the hypergeometric series.
        The solution to the hypergeometric equation is:
        \[
        y(x) = \sum_{n=0}^\infty a_n x^{n+s},
        \]
        where \(s\) is one of the roots of the indicial equation, and \(a_n\) are determined by the recurrence relation.
        The series converges for \(|x| < 1\).

\section*{Textbook question 7.6.3}
    Using the Wronskian determinant, show that the set of functions 
    \[
        \left(1, \frac{x^n}{n!} (n=1, 2, ..., N) \right)
    \]
    is linearly independent

    \subsection*{Solution:}
        The Wronskian determinant of a set of \(n\) functions \(f_1(x), f_2(x), \ldots, f_n(x)\) is defined as:
        \[
        W(f_1, f_2, \ldots, f_n)(x) = 
        \begin{vmatrix}
        f_1(x) & f_2(x) & \cdots & f_n(x) \\
        f_1'(x) & f_2'(x) & \cdots & f_n'(x) \\
        \vdots & \vdots & \ddots & \vdots \\
        f_1^{(n-1)}(x) & f_2^{(n-1)}(x) & \cdots & f_n^{(n-1)}(x)
        \end{vmatrix}.
        \]
        If \(W(f_1, f_2, \ldots, f_n)(x) \neq 0\) for some \(x\), the functions \(f_1(x), f_2(x), \ldots, f_n(x)\) are linearly independent.
        The given functions are:
        \[
        f_1(x) = 1, \quad f_2(x) = \frac{x}{1!}, \quad f_3(x) = \frac{x^2}{2!}, \quad \ldots, \quad f_n(x) = \frac{x^{n-1}}{(n-1)!}.
        \]
        To construct the Wronskian matrix, we compute the derivatives of each function.
        \begin{enumerate}      
            \item For \(f_1(x) = 1\):
            \[
            f_1'(x) = 0, \quad f_1''(x) = 0, \quad \ldots
            \]
            \item For \(f_2(x) = \frac{x}{1!}\):
            \[
            f_2'(x) = 1, \quad f_2''(x) = 0, \quad \ldots
            \]
            \item For \(f_3(x) = \frac{x^2}{2!}\):
            \[
            f_3'(x) = \frac{2x}{2!} = \frac{x}{1!}, \quad f_3''(x) = \frac{2}{2!} = \frac{1}{1!}, \quad f_3^{(3)}(x) = 0, \quad \ldots
            \]
            \item In general, for \(f_k(x) = \frac{x^{k-1}}{(k-1)!}\):
            \[
            f_k^{(m)}(x) =
            \begin{cases} 
            \frac{(k-1)(k-2)\cdots(k-m)}{(k-1)!} x^{k-1-m} & \text{if } m < k-1, \\
            \frac{1}{(k-1)!} (k-1)! & \text{if } m = k-1, \\
            0 & \text{if } m > k-1.
            \end{cases}
            \]
        \end{enumerate}
        The Wronskian matrix for the set of functions \(\left\{1, \frac{x}{1!}, \frac{x^2}{2!}, \ldots, \frac{x^{n-1}}{(n-1)!}\right\}\) is:
        \[
        W(f_1, f_2, \ldots, f_n) =
        \begin{vmatrix}
        1 & \frac{x}{1!} & \frac{x^2}{2!} & \cdots & \frac{x^{n-1}}{(n-1)!} \\
        0 & 1 & \frac{x}{1!} & \cdots & \frac{x^{n-2}}{(n-2)!} \\
        0 & 0 & 1 & \cdots & \frac{x^{n-3}}{(n-3)!} \\
        \vdots & \vdots & \vdots & \ddots & \vdots \\
        0 & 0 & 0 & \cdots & 1
        \end{vmatrix}.
        \]
        The Wronskian matrix is a triangular matrix, where all entries below the diagonal are zero. The diagonal elements of the matrix are:
        \[
        1, \quad 1, \quad 1, \quad \ldots, \quad 1.
        \]
        The determinant of a triangular matrix is the product of its diagonal elements. Therefore:
        \[
        \det(W(f_1, f_2, \ldots, f_n)) = 1.
        \]
        Since \( \det(W(f_1, f_2, \ldots, f_n)) = 1 \neq 0 \), the Wronskian is non-zero. This proves that the set of functions:
        \[
        \left\{1, \frac{x}{1!}, \frac{x^2}{2!}, \ldots, \frac{x^{n-1}}{(n-1)!}\right\}
        \]
        is linearly independent.
    
\section*{Textbook Problem 7.7.4}
    Find the general solutions to the following inhomogeneous ODE's:
    \[
        y'' - 3y' + 2y = sin(x)
    \]

    \subsection*{Solution:}
        The homogeneous equation is:
        \[
        y_h'' - 3y_h' + 2y_h = 0.
        \]
        The characteristic equation is:
        \[
        r^2 - 3r + 2 = 0.
        \]
        Factorizing:
        \[
        (r - 1)(r - 2) = 0.
        \]
        Thus, the roots are:
        \[
        r = 1, \quad r = 2.
        \]
        The general solution to the homogeneous equation is:
        \[
        y_h(x) = C_1 e^x + C_2 e^{2x},
        \]
        where \(C_1\) and \(C_2\) are arbitrary constants.
        The non-homogeneous term is \(\sin x\). To find a particular solution, we use the method of undetermined coefficients. Since the right-hand side is \(\sin x\), we try:
        \[
        y_p = A \cos x + B \sin x,
        \]
        where \(A\) and \(B\) are constants to be determined.
        The derivatives of \(y_p\) are:
        \[
        y_p' = -A \sin x + B \cos x,
        \]
        \[
        y_p'' = -A \cos x - B \sin x.
        \]
        Substitute \(y_p\), \(y_p'\), and \(y_p''\) into the differential equation:
        \[
        (-A \cos x - B \sin x) - 3(-A \sin x + B \cos x) + 2(A \cos x + B \sin x) = \sin x.
        \]
        Simplify term by term:
        \[
        (-A \cos x - B \sin x) + 3A \sin x - 3B \cos x + 2A \cos x + 2B \sin x = \sin x.
        \]
        Group terms involving \(\cos x\) and \(\sin x\):
        \[
        (-A - 3B + 2A)\cos x + (-B + 3A + 2B)\sin x = \sin x.
        \]
        Simplify coefficients:
        \[
        (A - 3B)\cos x + (3A + B)\sin x = \sin x.
        \]
        Equating the coefficients of \(\cos x\) and \(\sin x\):
        \begin{align*}
        A - 3B &= 0, \tag{1} \\
        3A + B &= 1. \tag{2}
        \end{align*}
        From \((1)\), solve for \(A\):
        \[
        A = 3B.
        \]
        Substitute \(A = 3B\) into \((2)\):
        \[
        3(3B) + B = 1,
        \]
        \[
        9B + B = 1,
        \]
        \[
        10B = 1 \implies B = \frac{1}{10}.
        \]
        Substitute \(B = \frac{1}{10}\) into \(A = 3B\):
        \[
        A = 3 \cdot \frac{1}{10} = \frac{3}{10}.
        \]
        Thus, the particular solution is:
        \[
        y_p = \frac{3}{10} \cos x + \frac{1}{10} \sin x.
        \]
        The general solution is the sum of the homogeneous solution and the particular solution:
        \[
        y(x) = y_h(x) + y_p(x).
        \]
        Substitute \(y_h(x)\) and \(y_p(x)\):
        \[
        y(x) = C_1 e^x + C_2 e^{2x} + \frac{3}{10} \cos x + \frac{1}{10} \sin x.
        \]
        The general solution to the differential equation is:
        \[
        y(x) = C_1 e^x + C_2 e^{2x} + \frac{3}{10} \cos x + \frac{1}{10} \sin x,
        \]
        where \(C_1\) and \(C_2\) are arbitrary constants.
    
\section*{Textbook Problem 7.8.1}
    Consider the Riccati equation 
    \[
        y' = y^2 - y - 2
    \]
    A particular solution to this equation is $y=2$. Find a more general solution.

    \subsection*{Solution:}
        To solve the Riccati equation, we use the substitution:
        \[
        y = y_p + \frac{1}{v},
        \]
        where \(y_p = 2\) is the particular solution, and \(v\) is a new function to solve for.
        The derivative becomes:
        \[
        y' = y_p' - \frac{1}{v^2}v',
        \]
        where \(y_p' = 0\) since \(y_p = 2\) is constant. Thus:
        \[
        y' = -\frac{1}{v^2}v'.
        \]
        Substituting \(y = y_p + \frac{1}{v}\) and \(y' = -\frac{1}{v^2}v'\) into the Riccati equation:
        \[
        -\frac{1}{v^2}v' = \left(y_p + \frac{1}{v}\right)^2 - \left(y_p + \frac{1}{v}\right) - 2.
        \]
        Expand the right-hand side:
        \[
        \left(y_p + \frac{1}{v}\right)^2 = y_p^2 + \frac{2y_p}{v} + \frac{1}{v^2},
        \]
        and
        \[
        \left(y_p + \frac{1}{v}\right) = y_p + \frac{1}{v}.
        \]
        Substituting these:
        \[
        -\frac{1}{v^2}v' = \left(y_p^2 + \frac{2y_p}{v} + \frac{1}{v^2}\right) - \left(y_p + \frac{1}{v}\right) - 2.
        \]
        Simplify:
        \[
        -\frac{1}{v^2}v' = y_p^2 - y_p - 2 + \frac{2y_p}{v} - \frac{1}{v} + \frac{1}{v^2}.
        \]
        Since \(y_p = 2\) satisfies the quadratic \(y_p^2 - y_p - 2 = 0\), the equation reduces to:
        \[
        -\frac{1}{v^2}v' = \frac{2y_p}{v} - \frac{1}{v} + \frac{1}{v^2}.
        \]
        Substituting \(y_p = 2\):
        \[
        -\frac{1}{v^2}v' = \frac{4 - 1}{v} + \frac{1}{v^2}.
        \]
        \[
        -\frac{1}{v^2}v' = \frac{3}{v} + \frac{1}{v^2}.
        \]
        Multiply through by \(-v^2\):
        \[
        v' = -3v - 1.
        \]
        The transformed equation is:
        \[
        v' + 3v = -1.
        \]
        This is a first-order linear differential equation. The standard form is:
        \[
        v' + P(x)v = Q(x),
        \]
        where \(P(x) = 3\) and \(Q(x) = -1\).
        The integrating factor is:
        \[
        \mu(x) = e^{\int P(x)dx} = e^{\int 3dx} = e^{3x}.
        \]
        Multiply through by \(\mu(x) = e^{3x}\):
        \[
        e^{3x}v' + 3e^{3x}v = -e^{3x}.
        \]
        The left-hand side is the derivative of \(e^{3x}v\):
        \[
        \frac{d}{dx}(e^{3x}v) = -e^{3x}.
        \]
        Integrate both sides:
        \[
        e^{3x}v = \int -e^{3x} dx,
        \]
        \[
        e^{3x}v = -\frac{1}{3}e^{3x} + C,
        \]
        where \(C\) is the constant of integration.
        Divide through by \(e^{3x}\):
        \[
        v = -\frac{1}{3} + Ce^{-3x}.
        \]
        Recall the substitution:
        \[
        y = y_p + \frac{1}{v}.
        \]
        Substitute \(v = -\frac{1}{3} + Ce^{-3x}\):
        \[
        y = 2 + \frac{1}{-\frac{1}{3} + Ce^{-3x}}.
        \]
        Simplify the fraction:
        \[
        y = 2 + \frac{1}{-\frac{1}{3} + Ce^{-3x}} = 2 + \frac{1}{\frac{-1 + 3Ce^{-3x}}{3}} = 2 + \frac{3}{-1 + 3Ce^{-3x}}.
        \]
        The general solution to the Riccati equation is:
        \[
        y(x) = 2 + \frac{3}{-1 + 3Ce^{-3x}},
        \]
        where \(C\) is an arbitrary constant.

\end{document}
